\chapter*{Введение}
\addcontentsline{toc}{chapter}{Введение}

Малые и средние предприятия часто принимают управленческие решения на основе разрозненных Excel-выгрузок, банковских выписок и данных учетных систем. Даже там, где данные есть, они хранятся фрагментированно и не сведены в единый, регулярно обновляемый отчёт -- собственник видит общий оборот, но не имеет быстрого ответа на вопросы: какие категории расходов съедают прибыль, почему возникают кассовые разрывы, насколько бизнес зависит от отдельных контрагентов и какие операции требуют ручной проверки.

Тема данной курсовой работы -- применение методов машинного обучения и больших языковых моделей для финансового и операционного анализа малого производственного бизнеса, и разработка инструмента, который делает этот анализ регулярным, а не разовым. В качестве практического примера рассматривается компания, производящая кондитерские изделия и выпечку, продающая продукцию через сети дискаунтеров, магазины низких цен и торговые рынки Москвы, Московской и Тульской областей.

Актуальность работы в том, что готовые BI-решения и модули учёта крупных ERP-систем избыточны и дороги для предприятия такого масштаба, а типовые отчёты 1С и таблицы Excel не воспроизводят ни управленческий ДДС с разбивкой по классифицированным операциям, ни концентрацию контрагентов, ни прогноз спроса и возвратов. В результате накопленные транзакционные данные -- о продажах, закупках, возвратах, аренде, налогах, платежах поставщикам -- остаются бухгалтерским массивом строк, а не инструментом управления.

Цель работы -- разработать и проверить на данных реального предприятия подход к применению ML и LLM для анализа транзакционных и операционных данных малого производственного бизнеса и встроить его в собственную аналитическую систему, повышающую прозрачность денежных потоков и формирующую основу для управленческих решений.

Для достижения цели поставлены следующие задачи:
\begin{enumerate}
  \item описать специфику применения ML и LLM в малом и среднем бизнесе;
  \item провести обзор источников данных компании и построить воспроизводимый пайплайн их очистки и подготовки;
  \item разработать схему классификации банковских операций и сравнить rule-based и LLM-подходы к разметке;
  \item сформировать управленческий cash flow и провести финансовый аудит денежных потоков;
  \item обучить и оценить прогнозные модели спроса, возвратов и денежного потока;
  \item реализовать собственную CRM-платформу, объединяющую загрузку данных, классификацию операций и прогнозные модели в одном интерфейсе;
  \item подготовить рекомендации по внедрению результатов в управленческую практику.
\end{enumerate}

Объект исследования -- малое производственное предприятие в сфере кондитерских изделий. Предмет исследования -- методы автоматической обработки, классификации и прогнозирования транзакционных и операционных данных предприятия, а также архитектура инструмента, делающего эти методы доступными собственнику без привлечения сторонних BI-решений.

Практическая значимость работы состоит в том, что предложенный подход переводит разовую работу с выгрузками в воспроизводимый контур: данные очищаются, операции классифицируются, денежный поток пересчитывается, модели прогнозирования обновляются, спорные операции выделяются для ручного контроля -- и всё это доступно в одном месте, а не пересобирается вручную каждый раз заново.
